\begin{frame}{Проблема и граница вывода}
Большие синтетические выборки не заменяют независимых людей, а реакция Drosophila
на вредоносное воздействие не эквивалентна субъективной боли человека.
\begin{block}{Проверяемая идея}
Выяснить, улучшает ли биологически ограниченное синтетическое предобучение перенос
на независимые человеческие наборы относительно human-only и контрольной синтетики.
\end{block}
\Evidence{S029, S062, S731}
\end{frame}
\begin{frame}{Шесть раздельных сущностей}
\small
\begin{tabularx}{\linewidth}{lX}
\toprule Сущность & Допустимый вывод \\ \midrule
Ноксический стимул & параметры воздействия \\
Ноцицептивная реакция & измеримая реакция на воздействие \\
Защитное поведение & наблюдаемое поведение \\
Самоотчёт боли & оценка зарегистрированного отчёта человека \\
ECAP & вызванное рекрутирование нервных волокон \\
Ответ на SCS & заранее заданный клинический исход \\ \bottomrule
\end{tabularx}
\end{frame}
\begin{frame}{Три гипотезы}
\begin{enumerate}
  \item H1: воспроизводимая симуляция воздействия и реакции без утечки команды.
  \item H2: прирост на participant-disjoint человеческих данных относительно контролей.
  \item H3: отдельная проверка SCS только на patient-linked ECAP/telemetry и исходах.
\end{enumerate}
\Evidence{S090, S105, S280, S730}
\end{frame}
